\chapter{乘法}\label{ch:g3:mult}

三袋弹珠，每袋 7 颗：可以算 $7 + 7 + 7$，但还有
一种更快的方法——乘法，即重复加法的运算。
本章先借助图形建立乘法表，再
学习竖式乘法。

\section{乘法的含义}

\begin{definition}[乘法]\label{def:g3:mult:def}
$3 \times 7$ 表示“三个七”：
\[
3 \times 7 = 7 + 7 + 7 = 21 .
\]
在图上，$3 \times 7$ 是一个点阵长方形，共 $3$ 行，
每行 $7$ 个圆点。
\end{definition}

\begin{omfigure}
\begin{tikzpicture}[scale=0.42]
  \foreach \r in {0,1,2}
    \foreach \c in {0,...,6}
      \fill[omDef] (\c,-\r) circle (0.27);
  \node[right, font=\small] at (7.0,-1.0) {$3$ 行，每行 $7$：\ \
  $3 \times 7 = 21$};
  \begin{scope}[yshift=-4.7cm]
  \foreach \r in {0,...,6}
    \foreach \c in {0,1,2}
      \fill[omThm] (\c,-\r) circle (0.27);
  \node[right, font=\small] at (7.0,-3.0) {$7$ 行，每行 $3$：\ \
  $7 \times 3 = 21$ 也成立！};
  \end{scope}
\end{tikzpicture}

{\small 同样的 $21$ 个圆点，无论按行数还是按列数，乘积的
次序不影响结果。知道 $3 \times 7$ 便能
直接得到 $7 \times 3$——要记的口诀少了一半！}
\end{omfigure}

\begin{proposition}[有用的规律]\label{prop:g3:mult:rules}
\begin{enumerate}
  \item 次序可以交换：$a \times b = b \times a$。
  \item 乘以 $1$ 不改变数值；乘以 $0$ 得
        $0$。
  \item 乘以 $10$ 就在末尾添一个零：$34 \times 10 = 340$；
        乘以 $100$ 则添两个零。
  \item 乘法表可以由相邻的结果推出：
        $6 \times 8 = 5 \times 8 + 8 = 40 + 8 = 48$。
\end{enumerate}
\end{proposition}
\admitted

\begin{example}[乘法表]\label{ex:g3:mult:table}
所有乘法口诀都排在同一张正方形表中。查 $6 \times 7$ 时：
先看 $6$ 行，再看 $7$ 列。
\begin{center}
\renewcommand{\arraystretch}{1.15}
\begin{tabular}{c|ccccccccc}
$\times$ & $1$ & $2$ & $3$ & $4$ & $5$ & $6$ & $7$ & $8$ & $9$ \\
\hline
$1$ & 1 & 2 & 3 & 4 & 5 & 6 & 7 & 8 & 9 \\
$2$ & 2 & 4 & 6 & 8 & 10 & 12 & 14 & 16 & 18 \\
$3$ & 3 & 6 & 9 & 12 & 15 & 18 & 21 & 24 & 27 \\
$4$ & 4 & 8 & 12 & 16 & 20 & 24 & 28 & 32 & 36 \\
$5$ & 5 & 10 & 15 & 20 & 25 & 30 & 35 & 40 & 45 \\
$6$ & 6 & 12 & 18 & 24 & 30 & 36 & 42 & 48 & 54 \\
$7$ & 7 & 14 & 21 & 28 & 35 & 42 & 49 & 56 & 63 \\
$8$ & 8 & 16 & 24 & 32 & 40 & 48 & 56 & 64 & 72 \\
$9$ & 9 & 18 & 27 & 36 & 45 & 54 & 63 & 72 & 81 \\
\end{tabular}
\end{center}
这张表关于对角线对称——这正是规律 1
的直观体现。
\end{example}

\section{竖式乘法}

\begin{method}[乘以一位数]\label{met:g3:mult:column}
计算 $234 \times 6$：
\begin{enumerate}
  \item 先算个位：$4 \times 6 = 24$，写 $4$，进 $2$；
  \item 再算十位，并加上进位：
        $3 \times 6 + 2 = 20$，写 $0$，进 $2$；
  \item 最后算百位，并加上进位：
        $2 \times 6 + 2 = 14$，写 $14$。
\end{enumerate}
\[
\begin{array}{r}
2\,3\,4 \\
\times \quad 6 \\
\hline
1\,4\,0\,4
\end{array}
\]
用估算检验：$234$ 接近 $200$，而
$200 \times 6 = 1\,200$，所以结果 $1\,404$ 是合理的。
\end{method}

\begin{example}[拆分后心算]\label{ex:g3:mult:split}
竖式乘法实际上是按数位拆分数。心算时也可以
用同样的办法：
\[
23 \times 4 = (20 + 3) \times 4 = 20 \times 4 + 3 \times 4
= 80 + 12 = 92 .
\]
还可借助一个好算的相近数：$99 \times 5 = 100 \times 5 - 5 =
495$。
\end{example}

\begin{omfigure}
\begin{tikzpicture}[scale=0.55]
  \fill[omDef!18] (0,0) rectangle (10,4);
  \fill[omThm!18] (10,0) rectangle (11.5,4);
  \draw[thick] (0,0) rectangle (11.5,4);
  \draw[thick] (10,0) -- (10,4);
  \node[font=\small] at (5,2) {$20 \times 4 = 80$};
  \node[font=\small, rotate=90] at (10.75,2) {$3 {\times} 4$};
  \node[below, font=\small] at (5,0) {$20$};
  \node[below, font=\small] at (10.75,0) {$3$};
  \node[left, font=\small] at (0,2) {$4$};
\end{tikzpicture}

{\small 为什么拆分有效：一个 $23 \times 4$ 的单位
方格长方形，可分为 $20 \times 4$ 的部分和 $3 \times 4$ 的部分——
共 $80 + 12 = 92$ 个方格。}
\end{omfigure}

\section{解决问题中的乘法}

\begin{example}\label{ex:g3:mult:problem}
一辆小巴可以载 $9$ 名乘客，$27$ 辆
小巴一共可以载多少名乘客？
\begin{enumerate}
  \item 这是乘法：$27$ 个 $9$。
  \item $27 \times 9 = 243$（竖式，或
        $27 \times 9 = 27 \times 10 - 27 = 270 - 27$）。
  \item 答：小巴一共可以载 $243$ 名乘客。
\end{enumerate}
\end{example}

\section{练习}

\begin{exercise}[$\star$]\label{exo:g3:mult:1}
写成乘法算式，再计算：
$8 + 8 + 8 + 8$；\quad $5 + 5 + 5$；\quad
$20 + 20 + 20 + 20 + 20$。
\end{exercise}

\begin{exercise}[$\star$]\label{exo:g3:mult:2}
背诵并填出：$6 \times 7$；\ $8 \times 4$；\ $9 \times 9$；
\ $7 \times 8$；\ $6 \times 6$。
\end{exercise}

\begin{exercise}[$\star$]\label{exo:g3:mult:3}
分别画出 $4 \times 6$ 的点阵长方形，再画出 $6 \times 4$ 的点阵长方形。
你发现这两个点数之间有什么联系？
\end{exercise}

\begin{exercise}[$\star$]\label{exo:g3:mult:4}
利用\cref{prop:g3:mult:rules}中的规律计算：
$56 \times 10$；\quad $8 \times 100$；\quad $73 \times 0$；\quad
$45 \times 1$；\quad $30 \times 10$。
\end{exercise}

\begin{exercise}[$\star$]\label{exo:g3:mult:5}
用竖式计算：$132 \times 3$；\quad $217 \times 4$；\quad
$408 \times 7$。
\end{exercise}

\begin{exercise}[$\star$]\label{exo:g3:mult:6}
用拆分的方法口算（见\cref{ex:g3:mult:split}）：
$31 \times 5$；\quad $42 \times 3$；\quad $25 \times 6$。
\end{exercise}

\begin{exercise}[$\star$]\label{exo:g3:mult:7}
借助好算的邻近数计算：$99 \times 7$；\quad
$101 \times 6$；\quad $98 \times 5$。
\end{exercise}

\begin{exercise}[$\star$]\label{exo:g3:mult:8}
先估算，再计算：$389 \times 5$（用 $400$ 估算）；
$612 \times 8$（用 $600$ 估算）。
\end{exercise}

\begin{exercise}[$\star$]\label{exo:g3:mult:9}
一个鸡蛋盒可以装 $6$ 个鸡蛋，$38$ 盒一共有多少个鸡蛋？请列出
算式、计算，并用一句话作答。
\end{exercise}

\begin{exercise}[$\star\star$]\label{exo:g3:mult:10}
一间教室有 $8$ 排，每排 $4$ 张课桌；每张课桌坐 $2$ 名学生。
这间教室可以坐多少名学生？（需要用两次乘法计算。）
\end{exercise}

\begin{exercise}[$\star\star$]\label{exo:g3:mult:11}
把 $24$ 把椅子摆成若干排，每排数量相同，求出所有排法（每排
$1$ 把、$2$ 把、\dots）。哪些排法可行？（利用乘法口诀表，在表中找
$24$ 的位置。）
\end{exercise}

\begin{exercise}[$\star\star$]\label{exo:g3:mult:12}
佐伊这样计算 $47 \times 6$：$40 \times 6 = 240$，
$7 \times 6 = 42$，再算 $240 + 42 = 282$。莉亚这样计算
$47 \times 6 = 50 \times 6 - 3 \times 6 = 300 - 18 = 282$。请解释
每种方法的思路。你更喜欢哪一种方法计算 $38 \times 4$？请用
两种方法都算一算。
\end{exercise}
